
\section{$A_\infty$ 双模的同伦等价：几何实现}
\label{sec:5-3}
\liftIndex{I193}\liftIndex{I194}

\subsection{滤过 $A_\infty$ 双模同态的构造}
\label{subsec:5-3-1}

本小节要证明：
由一对相对 spin 的拉格朗日子流形
构造出的滤过 $A_\infty$ 双模，
其同伦型不依赖于构造时所作的各种选择。
定理 5.3.1 是本小节的主要结果。

设 $(L^{(0)},L^{(1)})$ 为
$M$ 中一对相对 spin 的拉格朗日子流形，
$\psi\colon M\to M'$ 为辛微分同胚。
令 $\{J_t\}_t$、$\{J'_t\}_t$
（$t\in[0,1]$）
分别为 $M$、$M'$ 上的相容殆复结构族。
在 $i=0,1$ 处，
于 $M'$ 上取殆复结构同伦
$\{J^{(i)}_\rho\}_\rho$，
从 $J_i^\psi=\psi_*(J_i)$ 连接到 $J'_i$：
\begin{center}
\[
\begin{array}{ccc}
J^\psi_1=J^{(1)}_0
 &\xrightarrow{\;J^{(1)}_\rho\;}&
J^{(1)}_1=J'_1\\[2mm]
\Big\updownarrow\vcenter{\hbox{\scriptsize$J^\psi_t$}}
 &&\Big\updownarrow\vcenter{\hbox{\scriptsize$J'_t$}}\\[2mm]
J^\psi_0=J^{(0)}_0
 &\xrightarrow{\;J^{(0)}_\rho\;}&
J^{(0)}_1=J'_0 .
\end{array}
\]
\renewcommand{\thefigure}{5.3.1}
\refstepcounter{figure}
{\normalfont\bfseries 图式~\thefigure}\label{diag:5-3-1}
\end{center}

取多值截面
$\mathfrak s_i^{\operatorname{top}(\rho)}$
（$i=0,1$）。
第 4.6 节由此定义滤过
$A_\infty$ 同态
\[
\begin{aligned}
\mathfrak f^{i,\psi,\{J^{(i)}_\rho\}_\rho,
  \mathfrak s_i^{\operatorname{top}(\rho)}}:
&(C_1(L^{(i)};\Lambda_{0,\mathrm{nov}}),
  \mathfrak m_k^{J_i,\mathfrak s_i})
\\
&\longrightarrow
(C_3(L^{(i)\prime};\Lambda_{0,\mathrm{nov}}),
  \mathfrak m_k^{J'_i,\mathfrak s'_i}),
\\[1mm]
\mathfrak g^{i,\psi,\{J^{(i)}_\rho\}_\rho,
  \mathfrak s_i^{\operatorname{top}(\rho)}}:
&(C_1(L^{(i)};\Lambda_{0,\mathrm{nov}}),
  \mathfrak m_k^{J_i,\mathfrak s_i})
\\
&\longrightarrow
(C_2(L^{(i)\prime};\Lambda_{0,\mathrm{nov}}),
  \mathfrak m_k^{J'_i,\mathfrak s'_i}).
\end{aligned}
\]
以下将这两个带完整上标的同态
简记为 $\mathfrak f_i$、$\mathfrak g_i$。
这里 $L^{(i)\prime}=\psi(L^{(i)})$。
$C_1,C_2$ 分别由
$L^{(i)},L^{(i)\prime}$
中光滑奇异单形的可数集
$\mathcal X_1(L^{(i)})$、
$\mathcal X_2(L^{(i)\prime})$
确定；这些可数集在第 3.5 节
（及第 7.2 节）选定，
用以构造两侧的滤过 $A_\infty$ 代数结构。
$C_3$ 则由另一链的可数集
$\mathcal X_3(L^{(i)\prime})
 \supset\mathcal X_2(L^{(i)\prime})$
确定；这是在定理 4.6.1 中为构造
$\mathfrak f_i$ 所作的选择。
$\mathfrak g_i$ 定义为
$\mathfrak f_i$ 与下列自然包含诱导的
同态之同伦逆的复合：
\[
(C_2(L^{(i)\prime};\Lambda_{0,\mathrm{nov}}),
 \mathfrak m_k^{J'_i,\mathfrak s'_i})
\longrightarrow
(C_3(L^{(i)\prime};\Lambda_{0,\mathrm{nov}}),
 \mathfrak m_k^{J'_i,\mathfrak s'_i}),
\]
其中的包含是
$\mathcal X_2(L^{(i)\prime})
\subset\mathcal X_3(L^{(i)\prime})$。
参见第 4.6.1 小节开头的定理 4.6.1。

先考察 $L^{(0)}\cap L^{(1)}$ 与
$L^{(0)\prime}\cap L^{(1)\prime}$
均横截的情形，
并研究第 3.7.4 小节定义的
\[
C(L^{(1)},L^{(0)};\Lambda_{0,\mathrm{nov}}),
\qquad
C(L^{(1)\prime},L^{(0)\prime};
  \Lambda_{0,\mathrm{nov}}).
\]
这一讨论可以推广到
$L^{(0)}$、$L^{(1)}$ 干净相交的
博特--莫尔斯情形。
我们将在由两个不同的哈密顿同痕
分别移动两个拉格朗日子流形时
讨论该情形；
只须把那里的论证适当改用于此处。

本小节的主要结果如下。
\begin{theorem}{5.3.1}\label{thm:5-3-1}
存在滤过 $A_\infty$ 双模同态
\[
\varphi\colon
C(L^{(1)},L^{(0)};\Lambda_{0,\mathrm{nov}})
\longrightarrow
C(L^{(1)\prime},L^{(0)\prime};
 \Lambda_{0,\mathrm{nov}})
\]
立于 $(\mathfrak g_1,\mathfrak g_0)$ 之上，
且它是同伦等价。
\end{theorem}
注意，在定理 5.3.1 中，
两个拉格朗日子流形由同一个辛微分同胚
$\psi$ 同时移动。
此时以 $\Lambda_{0,\mathrm{nov}}$ 为系数的
弗洛尔上同调有良好定义。
在第 5.3.2--5.3.4 小节，
我们将研究由不同哈密顿同痕
分别移动两个拉格朗日子流形的
较微妙情形。
在那里仍能构造同伦等价，
但须以 $\Lambda_{\mathrm{nov}}$ 为系数环，
且所得只是第 5.2.1 小节引入的
弱滤过 $A_\infty$ 双模同态。

\begin{proof}
证明类似于定理 4.6.1，
但这里需要一个相容殆复结构的双参数族。
先注意，可选取 $J^{(i)}_\rho$，
使 $\rho<\epsilon$ 时
$J^{(i)}_\rho=J^{(i)}_0$，
而 $1-\epsilon<\rho$ 时
$J^{(i)}_\rho=J^{(i)}_1$；
这里 $\epsilon>0$ 足够小。
于是存在微分同胚
$\chi\colon\mathbb R\to(0,1)$，
满足\translatornote{原书实体页 306 的
(5.3.2.2) 写成未定义的
$J^\psi_{\chi(\tau)}$
从 $J^\psi_0$ 变到 $J^\psi_1$；
根据紧前所取的两族
$J^{(i)}_\rho$ 及 (5.3.3.1)，
此处按各自端点常值条件排印。
见疑误 A-0293。}
\begin{align}
\chi(\tau_1)&<\chi(\tau_2)
\quad(-A\le\tau_1<\tau_2\le A),
\tag{5.3.2.1}\label{eq:5-3-2-1}\\
J^{(i)}_{\chi(\tau)}
&=J^{(i)}_0\quad(\tau<-A),
\qquad
J^{(i)}_{\chi(\tau)}
=J^{(i)}_1\quad(\tau>A)
\quad(i=0,1).
\tag{5.3.2.2}\label{eq:5-3-2-2}
\end{align}
取 $M'$ 上的相容殆复结构双参数族
$\{J_{\tau,t}\}_{\tau,t}$，
$\tau\in\mathbb R$、$t\in[0,1]$，
满足
\begin{align}
J_{\tau,0}&=J^{(0)}_{\chi(\tau)},
\qquad
J_{\tau,1}=J^{(1)}_{\chi(\tau)},
\tag{5.3.3.1}\label{eq:5-3-3-1}\\
J_{\tau,t}
&=J^\psi_t\quad(\tau<-A),
\qquad
J_{\tau,t}=J'_t\quad(\tau>A).
\tag{5.3.3.2}\label{eq:5-3-3-2}
\end{align}
回忆 $J^{(i)}_0=J_i^\psi$。
这样的双参数族存在，
因为相容殆复结构的全体可缩。

令
$[\ell_p,w]\in Cr(L^{(0)},L^{(1)})$，
$[\ell_q,w']\in
 Cr(L^{(0)\prime},L^{(1)\prime})$；
又记
\[
\psi([\ell_p,w])
=[\psi\circ\ell_p,\psi\circ w]
\in Cr(L^{(0)\prime},L^{(1)\prime}).
\]
现在定义模空间\translatornote{原书实体页 306
在首次写该模空间时，
负端参数未加 $\psi$，
但其条带取值于 $M'$、
负端极限是 $\psi(p)$，
且实体页 307 的评价映射
已把同一参数写成
$\psi([\ell_p,w])$。
这里统一采用后者；见疑误 A-0294。}
\[
\mathcal M^\circ_{k_1,k_0}
\bigl(L^{(1)\prime},L^{(0)\prime};
 \{J_{\tau,t}\}_{\tau,t};
 \psi([\ell_p,w]),[\ell_q,w'];
 \operatorname{top}(\vec\tau)\bigr).
\]
如下给出定义。
令 $\Sigma$ 为
$\mathbb R\times[0,1]$
以及可能附着的球面气泡，
与前文相同。
考察满足以下条件的
$(\Sigma;u;\vec\tau^{(1)},\vec\tau^{(0)})$：
\begin{enumerate}[
 label=\textup{(5.3.4.\arabic*)},
 leftmargin=5.9em]
\item\label{cond:5-3-4-1}
$u\colon\Sigma\to M'$ 在条带上满足
\[
\frac{\partial u}{\partial\tau}
+J_{\tau,t}\frac{\partial u}{\partial t}=0.
\]
若在 $(\tau,t)$ 处附着球面气泡树，
该树中的每个球面均为
$J_{\tau,t}$-全纯的。
另要求稳定性条件。
\item\label{cond:5-3-4-2}
$u(\tau,0)\subset L^{(0)\prime}$，
$u(\tau,1)\subset L^{(1)\prime}$。
\item\label{cond:5-3-4-3}
$\lim_{\tau\to-\infty}u=\psi(p)$，
$\lim_{\tau\to+\infty}u=q$。
\item\label{cond:5-3-4-4}
\[
\begin{aligned}
\vec\tau^{(0)}
&=(\tau_1^{(0)},\ldots,\tau_{k_0}^{(0)}),
&-\infty<\tau_1^{(0)}<\cdots
 <\tau_{k_0}^{(0)}<+\infty,\\
\vec\tau^{(1)}
&=(\tau_1^{(1)},\ldots,\tau_{k_1}^{(1)}),
&-\infty<\tau_{k_1}^{(1)}<\cdots
 <\tau_1^{(1)}<+\infty.
\end{aligned}
\]
\end{enumerate}

\begin{figure}[htbp]
\centering
\setlength{\unitlength}{0.75mm}
\begin{picture}(104,38)
 \put(8,24){\line(1,0){88}}
 \put(8,12){\line(1,0){88}}
 \put(17,12){\line(1,2){6}}
 \put(29,12){\line(1,2){6}}
 \put(41,12){\line(1,2){6}}
 \put(53,12){\line(1,2){6}}
 \put(65,12){\line(1,2){6}}
 \put(77,12){\line(1,2){6}}
 \put(21,24){\makebox(0,0){$\times$}}
 \put(82,24){\makebox(0,0){$\times$}}
 \put(18,12){\makebox(0,0){$\times$}}
 \put(31,12){\makebox(0,0){$\times$}}
 \put(84,12){\makebox(0,0){$\times$}}
 \put(13,29){$\tau_{k_1}^{(1)}$}
 \put(77,29){$\tau_1^{(1)}$}
 \put(13,2){$\tau_1^{(0)}$}
 \put(28,2){$\tau_2^{(0)}$}
 \put(79,2){$\tau_{k_0}^{(0)}$}
 \put(46,3){$\cdots$}
 \put(46,28){$\cdots$}
\end{picture}
\par\smallskip
\renewcommand{\thefigure}{5.3.1}
\refstepcounter{figure}
{\normalfont\bfseries 图~\thefigure}\label{fig:5-3-1}
\end{figure}

上述模空间由所有这类
$(\Sigma;u;\vec\tau^{(0)},\vec\tau^{(1)})$
的 $\sim$ 等价类构成；
等价关系与定义 2.1.24 中的类似。
下文若未另作说明，
写 $\mathcal M^\circ_{k_1,k_0}$
时省去这里已经固定的
$L^{(i)\prime}$、$J_{\tau,t}$、
两端渐近类及 $\operatorname{top}(\vec\tau)$
等参数。
定义评价映射
\[
\begin{aligned}
\operatorname{ev}
&=(\operatorname{ev}^{(1)}_1,\ldots,
 \operatorname{ev}^{(1)}_{k_1};
 \operatorname{ev}^{(0)}_1,\ldots,
 \operatorname{ev}^{(0)}_{k_0})
\\
&\colon
\mathcal M^\circ_{k_1,k_0}
\longrightarrow
(L^{(1)})^{k_1}\times(L^{(0)})^{k_0}
\end{aligned}
\]
如下：
\[
\begin{aligned}
&\operatorname{ev}
(\Sigma;u;\vec\tau^{(1)},\vec\tau^{(0)})
\\
&\quad=
\bigl(
 \psi^{-1}(u(\tau^{(1)}_1,1)),\ldots,
 \psi^{-1}(u(\tau^{(1)}_{k_1},1));
\\
&\qquad
 \psi^{-1}(u(\tau^{(0)}_1,0)),\ldots,
 \psi^{-1}(u(\tau^{(0)}_{k_0},0))
\bigr).
\end{aligned}
\]
另定义时间评价映射
\[
\begin{aligned}
\operatorname{ev}_{\mathrm{time},(0)}
&=(\operatorname{ev}_{\mathrm{time},(0),1},
  \ldots,
  \operatorname{ev}_{\mathrm{time},(0),k_0})
\\
&\colon\mathcal M^\circ_{k_1,k_0}
\longrightarrow\mathbb R^{k_0},
\end{aligned}
\]
其中
\[
\operatorname{ev}_{\mathrm{time},(0),j}
(\Sigma;u;\vec\tau^{(1)},\vec\tau^{(0)})
=\tau_j^{(0)}.
\]
以同样方式定义
\[
\begin{aligned}
\operatorname{ev}_{\mathrm{time},(1)}
&=(\operatorname{ev}_{\mathrm{time},(1),1},
  \ldots,
  \operatorname{ev}_{\mathrm{time},(1),k_1})
\\
&\colon\mathcal M^\circ_{k_1,k_0}
\longrightarrow\mathbb R^{k_1}.
\end{aligned}
\]

为定义滤过 $A_\infty$ 双模同态，
须将上述模空间与第 4.6 节使用的
\[
\mathcal M_{k+1}^{\mathrm{main}}
\bigl(M',L^{(i)\prime},
 \{J^{(i)}_\rho\}_\rho;
 \operatorname{top}(\rho)\bigr)
\]
按下文方式结合。
注意参数 $\rho,\tau$ 满足
如下坐标变换：
\[
\rho=\chi(\tau),
\qquad \tau\in\mathbb R,\quad \rho\in[0,1].
\]
严格地说，
所取 $\chi\colon\mathbb R\to(0,1)$
是满足 (5.3.2) 的微分同胚，
而在 $\tau<-A$ 与 $\tau>A$ 时，
由 $\chi(\tau)$ 参数化的殆复结构
保持常值。

取非负整数
$l_1^{(0)},\ldots,l_a^{(0)}$
与 $l_1^{(1)},\ldots,l_b^{(1)}$，
使
\[
\sum_{i=1}^{a}l_i^{(0)}=k_0,
\qquad
\sum_{i'=1}^{b}l_{i'}^{(1)}=k_1.
\]
考虑系统
\[
\begin{aligned}
\mathfrak D
=\bigl(
 &(\Sigma;u;\vec\tau^{(1)},\vec\tau^{(0)});
\\[-1mm]
 &((\Sigma_{(1),i'},\vec z_{(1),i'}),
   (u_{(1),i',\alpha}),
   (\rho_{(1),i',\alpha}))_{i'=1}^{b};
\\[-1mm]
 &((\Sigma_{(0),i},\vec z_{(0),i}),
   (u_{(0),i,\alpha}),
   (\rho_{(0),i,\alpha}))_{i=1}^{a}
\bigr),
\end{aligned}
\]
并要求：\translatornote{原书实体页 308 的
(5.3.5.4)、(5.3.5.5)
把条带附着点写为
$u(\chi(\tau_i^{(p)}),p)$；
条带的第一坐标是 $\tau$，
而 $\rho=\chi(\tau)$ 是所附圆盘族的参数。
译文按实体页 307 的评价映射
将条带附着点写为
$u(\tau_i^{(p)},p)$，
仍在 (5.3.5.6) 用 $\chi(\tau)$
比较圆盘时间；见疑误 A-0295。}
\begin{enumerate}[
 label=\textup{(5.3.5.\arabic*)},
 leftmargin=5.9em]
\item\label{cond:5-3-5-1}
$(\Sigma;u;\vec\tau^{(1)},\vec\tau^{(0)})$
属于
\[
\mathcal M^\circ_{b,a}
\bigl(
L^{(1)\prime},L^{(0)\prime};
\{J_{\tau,t}\}_{\tau,t};
\psi([\ell_p,w]),[\ell_q,w'];
\operatorname{top}(\vec\tau)
\bigr).
\]
\item\label{cond:5-3-5-2}
对每个 $i=1,\ldots,a$，
\[
\begin{aligned}
&((\Sigma_{(0),i},\vec z_{(0),i}),
 (u_{(0),i,\alpha}),
 (\rho_{(0),i,\alpha}))
\\[-1mm]
&\quad\in
\mathcal M_{l_i^{(0)}+1}^{\mathrm{main}}
\bigl(
 M';L^{(0)\prime},
 \{J^{(0)}_\rho\}_\rho:\beta_i^{(0)};
 \operatorname{top}(\rho)
\bigr).
\end{aligned}
\]
\item\label{cond:5-3-5-3}
对每个 $i'=1,\ldots,b$，
\[
\begin{aligned}
&((\Sigma_{(1),i'},\vec z_{(1),i'}),
 (u_{(1),i',\alpha}),
 (\rho_{(1),i',\alpha}))
\\[-1mm]
&\quad\in
\mathcal M_{l_{i'}^{(1)}+1}^{\mathrm{main}}
\bigl(
 M';L^{(1)\prime},
 \{J^{(1)}_\rho\}_\rho:\beta_{i'}^{(1)};
 \operatorname{top}(\rho)
\bigr).
\end{aligned}
\]
\item\label{cond:5-3-5-4}
$\operatorname{ev}_0
((\Sigma_{(0),i},\vec z_{(0),i}),
 (u_{(0),i,\alpha}),
 (\rho_{(0),i,\alpha}))
=u(\tau_i^{(0)},0)$。
\item\label{cond:5-3-5-5}
$\operatorname{ev}_0
((\Sigma_{(1),i'},\vec z_{(1),i'}),
 (u_{(1),i',\alpha}),
 (\rho_{(1),i',\alpha}))
=u(\tau_{i'}^{(1)},1)$。
\item\label{cond:5-3-5-6}
$\rho_{(0),i,0}\le\chi(\tau_i^{(0)})$，
$\rho_{(1),i',0}\le\chi(\tau_{i'}^{(1)})$。
\item\label{cond:5-3-5-7}
$(\sum_i\beta_i^{(0)}
 +\sum_{i'}\beta_{i'}^{(1)})\#w
\sim w'$。
\end{enumerate}
回忆 $\rho_{(0),i,0}$
是 $D_{(0),i}$ 最后一个分支所对应的
“时间”；亦即
$u_{(0),i,0}$
是 $J^{(0)}_{\rho_{(0),i,0}}$-全纯的。
又回忆 $\chi\colon\mathbb R\to(0,1)$
已在 (5.3.2) 给出。

\begin{figure}[!ht]
\centering
\setlength{\unitlength}{0.8mm}
\begin{picture}(150,139)
 \put(8,91){\line(1,0){134}}
 \put(8,69){\line(1,0){134}}
 \put(75,105){\circle{28}}
 \put(75,128){\circle{18}}
 \put(75,91){\circle*{1}}
 \put(75,119){\circle*{1}}
 \put(45,55){\circle{28}}
 \put(30,36){\circle{16}}
 \put(60,36){\circle{16}}
 \put(45,69){\circle*{1}}
 \put(36,44){\circle*{1}}
 \put(54,44){\circle*{1}}
 \put(106,56){\circle{26}}
 \put(106,69){\circle*{1}}
 \put(26,91){\makebox(0,0){$\times$}}
 \put(61,105){\makebox(0,0){$\times$}}
 \put(89,105){\makebox(0,0){$\times$}}
 \put(129,91){\makebox(0,0){$\times$}}
 \put(24,31){\makebox(0,0){$\times$}}
 \put(66,31){\makebox(0,0){$\times$}}
 \put(73,69){\makebox(0,0){$\times$}}
 \put(129,69){\makebox(0,0){$\times$}}
 \put(26,95){\makebox(0,0){$\scriptstyle P^{(1)}_4$}}
 \put(52,109){\makebox(0,0){$\scriptstyle P^{(1)}_3$}}
 \put(99,109){\makebox(0,0){$\scriptstyle P^{(1)}_2$}}
 \put(129,95){\makebox(0,0){$\scriptstyle P^{(1)}_1$}}
 \put(17,28){\makebox(0,0){$\scriptstyle P^{(0)}_1$}}
 \put(74,28){\makebox(0,0){$\scriptstyle P^{(0)}_2$}}
 \put(73,73){\makebox(0,0){$\scriptstyle P^{(0)}_3$}}
 \put(129,73){\makebox(0,0){$\scriptstyle P^{(0)}_4$}}
 \put(75,128){\makebox(0,0){$\scriptscriptstyle\rho_{(1),2,\alpha_1}$}}
 \put(75,105){\makebox(0,0){$\scriptscriptstyle\rho_{(1),2,\alpha_0}$}}
 \put(45,55){\makebox(0,0){$\scriptscriptstyle\rho_{(0),1,\alpha_0}$}}
 \put(30,36){\makebox(0,0){$\scriptscriptstyle\rho_{(0),1,\alpha_2}$}}
 \put(60,36){\makebox(0,0){$\scriptscriptstyle\rho_{(0),1,\alpha_1}$}}
 \put(106,56){\makebox(0,0){$\scriptscriptstyle\rho_{(0),3,\alpha_0}$}}
 \put(26,83){\makebox(0,0){$\scriptstyle\tau^{(1)}_3$}}
 \put(75,83){\makebox(0,0){$\scriptstyle\tau^{(1)}_2$}}
 \put(129,83){\makebox(0,0){$\scriptstyle\tau^{(1)}_1$}}
 \put(45,74){\makebox(0,0){$\scriptstyle\tau^{(0)}_1$}}
 \put(73,62){\makebox(0,0){$\scriptstyle\tau^{(0)}_2$}}
 \put(106,74){\makebox(0,0){$\scriptstyle\tau^{(0)}_3$}}
 \put(129,62){\makebox(0,0){$\scriptstyle\tau^{(0)}_4$}}
 \put(5,8){$\scriptstyle
 \rho_{(0),1,\alpha_2},
 \rho_{(0),1,\alpha_1}
 \le\rho_{(0),1,\alpha_0}
 \le\chi(\tau^{(0)}_1)$}
 \put(80,8){$\scriptstyle
 \rho_{(0),3,\alpha_0}\le\chi(\tau^{(0)}_3)$}
 \put(80,1){$\scriptstyle
 \rho_{(1),2,\alpha_1}\le\rho_{(1),2,\alpha_0}
 \le\chi(\tau^{(1)}_2)$}
\end{picture}
\par\smallskip
\renewcommand{\thefigure}{5.3.2}
\refstepcounter{figure}
{\normalfont\bfseries 图~\thefigure}\label{fig:5-3-2}
\end{figure}

在记号的宽松用法下，
还允许 $l_j^{(i)}=1$、
$\beta_j^{(i)}=0$（$i=0$ 或 $1$）。
此时
$\mathcal M^{\mathrm{main}}_2
(M';L^{(i)\prime},
  \{J^{(i)}_\rho\}_\rho;
  0;\operatorname{top}(\rho))$
原先未定义；
现规定它等于单点集，
并且 (5.3.5.4)、(5.3.5.5)、
(5.3.5.6) 在该情形下均无要求。
见图 5.3.2。

令
\[
\mathcal M^{\mathrm{unbr}}_{k_1,k_0}
\bigl(
 L^{(1)\prime},L^{(0)\prime};
 \{J_{\tau,t}\}_{\tau,t};
 \psi([\ell_p,w]),[\ell_q,w'];
 \operatorname{top}(\vec\tau)
\bigr)
\]
为所有上述系统 $\mathfrak D$ 的集合。
它也可定义成这些模空间纤维积的并。
上标 $\mathrm{unbr}$ 意为“未断裂”。
定义评价映射
\[
\begin{aligned}
\operatorname{ev}
&=(\operatorname{ev}^{(1)}_1,\ldots,
  \operatorname{ev}^{(1)}_{k_1};
  \operatorname{ev}^{(0)}_1,\ldots,
  \operatorname{ev}^{(0)}_{k_0})
\\
&\colon
\mathcal M^{\mathrm{unbr}}_{k_1,k_0}
\longrightarrow
(L^{(1)})^{k_1}\times
(L^{(0)})^{k_0}.
\end{aligned}
\]
具体说明
$\operatorname{ev}^{(0)}_m$。
取 $i,j$ 使
\[
m=l_1^{(0)}+\cdots+l_{i-1}^{(0)}+j.
\]
则
\[
\begin{aligned}
\operatorname{ev}^{(0)}_m(\mathfrak D)
=\operatorname{ev}_j\bigl(
 &(\Sigma_{(0),i},\vec z_{(0),i}),
 (u_{(0),i,\alpha}),
 (\rho_{(0),i,\alpha})
\bigr),
\end{aligned}
\]
右端的 $\operatorname{ev}_j$
按第 4.6.1 小节的方式定义；
这对应于在第 $i$ 个气泡的
第 $j$ 个标记点取第 $m$ 个评价。
当 $l_i^{(0)}=1$、
$\beta_i^{(0)}=0$ 时，
规定
\[
\operatorname{ev}^{(0)}_m
=\operatorname{ev}_j
(\Sigma;u;\vec\tau^{(1)},\vec\tau^{(0)}),
\]
此时取的是
$\mathbb R\times\{0\}$
上的标记点，
其全局序号为
$m=l_1^{(0)}+\cdots+l_{i-1}^{(0)}+1$。\translatornote{原书实体页 309
把此特例的全局序号写成
$l_1^{(0)}+\cdots+
 l_{i-1}^{(0)}+1=j$，
但同段前文用 $m$ 标记全局序号，
用 $j$ 标记气泡内的局部序号。
这里沿用前文的 $m$；见疑误 A-0296。}
上边界的
$\operatorname{ev}^{(1)}_m$
以同样方式定义。

可将
$\mathcal M^{\mathrm{unbr}}_{k_1,k_0}$
紧化为
\[
\mathcal M_{k_1,k_0}
\bigl(
 L^{(1)\prime},L^{(0)\prime};
 \{J_{\tau,t}\}_{\tau,t};
 \psi([\ell_p,w]),[\ell_q,w'];
 \operatorname{top}(\vec\tau)
\bigr).
\]
这里定义
\[
\begin{aligned}
&\mathcal M_{k_1,k_0}
\bigl(
 L^{(1)\prime},L^{(0)\prime};
 \{J_{\tau,t}\}_{\tau,t};
 \psi([\ell_p,w]),[\ell_q,w'];
 \operatorname{top}(\vec\tau)
\bigr)
\\
&=
\coprod_{l_0,l_1}
\coprod_{m_1,m_0}
\coprod_{p',w''}
\coprod_{q',w'''}
\Bigl[
 \mathcal M_{l_1,l_0}
 \bigl(
 L^{(1)\prime},L^{(0)\prime};
 \{J_{0,t}\}_t;
 \psi([\ell_p,w]),[\ell_{p'},w'']
 \bigr)
\\
&\hspace{1.5em}
{}\times
 \mathcal M^{\mathrm{unbr}}_{
   k_1-l_1-m_1,\,
   k_0-l_0-m_0}
 \bigl(
 L^{(1)\prime},L^{(0)\prime};
 \{J_{\tau,t}\}_{\tau,t};
 [\ell_{p'},w''],[\ell_{q'},w'''];
 \operatorname{top}(\vec\tau)
 \bigr)
\\
&\hspace{1.5em}
{}\times
 \mathcal M_{m_1,m_0}
 \bigl(
 L^{(1)\prime},L^{(0)\prime};
 \{J_{1,t}\}_t;
 [\ell_{q'},w'''],[\ell_q,w']
 \bigr)
\Bigr].
\end{aligned}
\]
首尾两个因子按 (3.7.31) 定义；
写出 $\{J_{i,t}\}_t$，
是为标明所采用的单参数殆复结构族。

\begin{figure}[htbp]
\centering
\setlength{\unitlength}{0.85mm}
\begin{picture}(115,51)
 \qbezier(9,25)(24,47)(39,25)
 \qbezier(9,25)(24,3)(39,25)
 \qbezier(39,25)(54,47)(69,25)
 \qbezier(39,25)(54,3)(69,25)
 \qbezier(69,25)(84,47)(99,25)
 \qbezier(69,25)(84,3)(99,25)
 \put(2,18){\line(1,1){7}}
 \put(2,32){\line(1,-1){7}}
 \put(99,25){\line(1,1){7}}
 \put(99,25){\line(1,-1){7}}
 \put(9,25){\circle*{1}}
 \put(39,25){\circle*{1}}
 \put(69,25){\circle*{1}}
 \put(99,25){\circle*{1}}
 \put(17,37){\makebox(0,0){$\times$}}
 \put(24,40){\makebox(0,0){$\times$}}
 \put(31,37){\makebox(0,0){$\times$}}
 \put(17,13){\makebox(0,0){$\times$}}
 \put(24,10){\makebox(0,0){$\times$}}
 \put(31,13){\makebox(0,0){$\times$}}
 \put(45,34){\makebox(0,0){$\times$}}
 \put(53,39){\makebox(0,0){$\times$}}
 \put(61,34){\makebox(0,0){$\times$}}
 \put(45,16){\makebox(0,0){$\times$}}
 \put(53,11){\makebox(0,0){$\times$}}
 \put(61,16){\makebox(0,0){$\times$}}
 \put(76,36){\makebox(0,0){$\times$}}
 \put(84,39){\makebox(0,0){$\times$}}
 \put(92,36){\makebox(0,0){$\times$}}
 \put(76,14){\makebox(0,0){$\times$}}
 \put(84,11){\makebox(0,0){$\times$}}
 \put(92,14){\makebox(0,0){$\times$}}
 \put(13,24){\scriptsize$J_{0,t}$-全纯}
 \put(44,24){\scriptsize$J_{\tau,t}$-全纯}
 \put(73,24){\scriptsize$J_{1,t}$-全纯}
 \put(7,2){$p$}
 \put(38,2){$p'$}
 \put(68,2){$q'$}
 \put(98,2){$q$}
 \put(16,46){$l_0$}
 \put(47,46){$k_1-l_1-m_1$}
 \put(81,46){$m_0$}
 \put(16,1){$l_1$}
 \put(43,1){$k_0-l_0-m_0$}
 \put(82,1){$m_1$}
 \put(105,29){$L^{(1)}$}
 \put(105,18){$L^{(0)}$}
\end{picture}
\par\smallskip
\renewcommand{\thefigure}{5.3.3}
\refstepcounter{figure}
{\normalfont\bfseries 图~\thefigure}\label{fig:5-3-3}
\end{figure}

按与 (3.7.30)、(3.7.33) 类似的方式，
将 $\operatorname{ev}$
延拓到该紧化。
现令
$\vec P^{(i)}
=(P_1^{(i)},\ldots,P_{k_i}^{(i)})$，
其中
$P_j^{(i)}\in C_1(L^{(i)};\mathbb Q)$
（$i=0,1$），
定义纤维积
\[
\begin{aligned}
&\mathcal M_{k_1,k_0}
\bigl(
 L^{(1)\prime},L^{(0)\prime};
 \{J_{\tau,t}\}_{\tau,t};
 \psi([\ell_p,w]),[\ell_q,w'];
 \operatorname{top}(\vec\tau);
 \vec P^{(1)};\vec P^{(0)}
\bigr)
\\
&\quad=
\mathcal M_{k_1,k_0}
\bigl(
 L^{(1)\prime},L^{(0)\prime};
 \{J_{\tau,t}\}_{\tau,t};
 \psi([\ell_p,w]),[\ell_q,w'];
 \operatorname{top}(\vec\tau)
\bigr)
\\
&\hspace{2em}
{}_{\operatorname{ev}}\!\times
\left(
 \prod_jP_j^{(1)}
 \times\prod_iP_i^{(0)}
\right),
\end{aligned}
\]
纤维积的基是
$(L^{(1)})^{k_1}\times(L^{(0)})^{k_0}$。

\begin{proposition}{5.3.6}\label{prop:5-3-6}
\textup{(1)}
上述带 $\vec P^{(1)},\vec P^{(0)}$
的模空间具有定向仓西结构，
维数为
\begin{equation}
\mu([\ell_q,w'])
-\mu([\ell_p,w])
+k_0+k_1
-\sum_i\deg P_i^{(0)}
-\sum_i\deg P_i^{(1)}.
\tag{5.3.7}\label{eq:5-3-7}
\end{equation}
\textup{(2)}
其上存在多值截面 $\mathfrak s$，
在无穷远处与前已选取的
各多值截面相容。
记该模空间的零集为
\[
\mathcal M_{k_1,k_0}
\bigl(
 L^{(1)\prime},L^{(0)\prime};
 \{J_{\tau,t}\}_{\tau,t};
 \psi([\ell_p,w]),[\ell_q,w'];
 \operatorname{top}(\vec\tau);
 \vec P^{(1)};\vec P^{(0)}
\bigr)^{\mathfrak s}
=\mathfrak s^{-1}(0).
\]
它是维数由 (5.3.7) 给出的
$\mathbb Q$-链；
适当三角剖分后，
将其视为奇异链。
\end{proposition}
证明与推迟至第 7 章的其他陈述相同。

令
\[
T^{\lambda_j^{(i)}}e^{\mu_j^{(i)}}P_j^{(i)}
\in C_1(L^{(i)};\Lambda_{0,\mathrm{nov}})
\quad(i=0,1;\ j=1,\ldots,k_i),
\]
并令
\[
T^\lambda e^\mu[\ell_p,w]
\in C(L^{(1)},L^{(0)};\Lambda_{0,\mathrm{nov}}).
\]
定义
\[
\begin{aligned}
&\phi_{k_1,k_0}
\bigl(
 T^{\lambda_1^{(1)}}e^{\mu_1^{(1)}}P_1^{(1)},
 \ldots,
 T^{\lambda_{k_1}^{(1)}}e^{\mu_{k_1}^{(1)}}P_{k_1}^{(1)};
 T^\lambda e^\mu[\ell_p,w];
\\[-1mm]
&\hspace{9em}
 T^{\lambda_1^{(0)}}e^{\mu_1^{(0)}}P_1^{(0)},
 \ldots,
 T^{\lambda_{k_0}^{(0)}}e^{\mu_{k_0}^{(0)}}P_{k_0}^{(0)}
\bigr)
\\
&\qquad=
\sum_{[\ell_q,w']}
\#
\Bigl(
\\
&\qquad\quad
 \mathcal M_{k_1,k_0}
 \bigl(
 L^{(1)\prime},L^{(0)\prime};
 \{J_{\tau,t}\}_{\tau,t};
 \psi([\ell_p,w]),[\ell_q,w'];
 \operatorname{top}(\vec\tau);
 \vec P^{(1)};\vec P^{(0)}
 \bigr)^{\mathfrak s}
\Bigr)
\\
&\hspace{4em}
T^{\lambda'}e^{\mu'}[\ell_q,w'].
\end{aligned}
\]
求和遍历使所出现模空间
为零维的 $[\ell_q,w']$，并取
\[
\lambda'
=\sum_{i=1}^{k_0}\lambda_i^{(0)}
 +\lambda
 +\sum_{i=1}^{k_1}\lambda_i^{(1)},
\qquad
\mu'
=\sum_{i=1}^{k_0}\mu_i^{(0)}
 +\mu
 +\sum_{i=1}^{k_1}\mu_i^{(1)}.
\]

\begin{lemma}{5.3.8}\label{lem:5-3-8}
$\phi=\{\phi_{k_1,k_0}\}_{k_1,k_0}$
是立于滤过 $A_\infty$ 同态对
$(\mathfrak f_1,\mathfrak f_0)$ 之上的
滤过 $A_\infty$ 双模同态。
\end{lemma}
这里的 $\mathfrak f_0,\mathfrak f_1$
即前面写出完整上标的
第 4.6 节所定义同态。

\begin{proof}
为证明引理 5.3.8，
须研究虚维数为 $1$ 时
上述带 $\vec P^{(1)},\vec P^{(0)}$
的模空间之边界。
设 $\mathfrak D_j$（$j=1,2,\ldots$）
是模空间中的序列，
在其紧化中收敛到一个边界点。
可假定 $\mathfrak D_j$ 的组合类型
不随 $j$ 变化。写作
\[
\begin{aligned}
\mathfrak D_j
=\bigl(
 &(\Sigma_j;u_j;
   \vec\tau^{(1),(j)},\vec\tau^{(0),(j)});
\\
 &((\Sigma_{(1),i',(j)},\vec z_{(1),i',(j)}),
   (u_{(1),i',(j),\alpha}),
   (\rho_{(1),i',(j),\alpha}));
\\
 &((\Sigma_{(0),i,(j)},\vec z_{(0),i,(j)}),
   (u_{(0),i,(j),\alpha}),
   (\rho_{(0),i,(j),\alpha}))
\bigr).
\end{aligned}
\]
参见 (5.3.5)。
当 $j\to\infty$ 时，
出现下列情形之一：
\begin{description}[
 leftmargin=6em,
 labelwidth=5.8em,
 style=nextline]
\item[\textup{(5.3.9.1)}]
$u_j$ 分裂为两部分，
其中一部分滑向 $\tau\to-\infty$ 一端。
\item[\textup{(5.3.9.2)}]
$u_j$ 分裂为两部分，
其中一部分滑向 $\tau\to+\infty$ 一端。
\item[\textup{(5.3.9.3)}]
$u_j$ 在 $t=0$ 处产生气泡。
\item[\textup{(5.3.9.4)}]
$u_j$ 在 $t=1$ 处产生气泡。
\item[\textup{(5.3.9.5.$p$)}]
$\Sigma_{(p),i,(j)}$
的一个分支
$\Sigma_{(p),i,(j),\alpha}$
分裂为两部分，其中 $p=0,1$。
\item[\textup{(5.3.9.6.$p$)}]
\[
\lim_{j\to\infty}
\rho_{(p),i,(j),\alpha}
=
\lim_{j\to\infty}
\rho_{(p),i,(j),\alpha'},
\]
其中
$\Sigma_{(p),i,(j),\alpha}
\cap\Sigma_{(p),i,(j),\alpha'}$
非空。
\item[\textup{(5.3.9.7.$p$)}]
$\lim_{j\to\infty}
\rho_{(p),i,(j),\alpha}=0$。
\item[\textup{(5.3.9.8.$p$)}]
\[
\lim_{j\to\infty}
\rho_{(p),i,(j),0}
=
\lim_{j\to\infty}
\chi(\tau_i^{(p),(j)}).
\]
这里 $\alpha=0$ 是
包含第 $0$ 个标记点的分支之指标。
\end{description}
(5.3.9.8.$p$)
对应于 (5.3.5.6) 中
不等式取等号的情形。
参见图 5.3.4--5.3.10。

\begin{figure}[htbp]
\centering
\setlength{\unitlength}{0.9mm}
\begin{picture}(105,54)(0,-8)
 \qbezier(4,17)(25,43)(52,17)
 \qbezier(4,17)(25,-9)(52,17)
 \qbezier(52,17)(77,43)(101,17)
 \qbezier(52,17)(77,-9)(101,17)
 \put(18,27){\makebox(0,0){$\times$}}
 \put(27,30){\makebox(0,0){$\times$}}
 \put(37,26){\makebox(0,0){$\times$}}
 \put(19,7){\makebox(0,0){$\times$}}
 \put(28,4){\makebox(0,0){$\times$}}
 \put(38,8){\makebox(0,0){$\times$}}
 \put(66,27){\makebox(0,0){$\times$}}
 \put(76,30){\makebox(0,0){$\times$}}
 \put(87,27){\makebox(0,0){$\times$}}
 \put(66,7){\makebox(0,0){$\times$}}
 \put(76,4){\makebox(0,0){$\times$}}
 \put(87,7){\makebox(0,0){$\times$}}
 \put(17,17){\scriptsize$J_{0,t}$-全纯}
 \put(64,17){\scriptsize$J_{\tau,t}$-全纯}
 \put(98,22){$L^{(0)}$}
 \put(98,10){$L^{(1)}$}
\end{picture}
\par\smallskip
\renewcommand{\thefigure}{5.3.4}
\refstepcounter{figure}
{\normalfont\bfseries 图~\thefigure}
\quad\textup{对应 (5.3.9.1)}\label{fig:5-3-4}
\end{figure}

\begin{figure}[htbp]
\centering
\setlength{\unitlength}{0.9mm}
\begin{picture}(105,54)(0,-8)
 \qbezier(4,17)(29,43)(53,17)
 \qbezier(4,17)(29,-9)(53,17)
 \qbezier(53,17)(78,43)(101,17)
 \qbezier(53,17)(78,-9)(101,17)
 \put(18,27){\makebox(0,0){$\times$}}
 \put(29,30){\makebox(0,0){$\times$}}
 \put(39,26){\makebox(0,0){$\times$}}
 \put(18,7){\makebox(0,0){$\times$}}
 \put(29,4){\makebox(0,0){$\times$}}
 \put(39,8){\makebox(0,0){$\times$}}
 \put(67,27){\makebox(0,0){$\times$}}
 \put(78,30){\makebox(0,0){$\times$}}
 \put(89,27){\makebox(0,0){$\times$}}
 \put(67,7){\makebox(0,0){$\times$}}
 \put(78,4){\makebox(0,0){$\times$}}
 \put(89,7){\makebox(0,0){$\times$}}
 \put(16,17){\scriptsize$J_{\tau,t}$-全纯}
 \put(65,17){\scriptsize$J_{1,t}$-全纯}
 \put(98,22){$L^{(1)}$}
 \put(98,10){$L^{(0)}$}
\end{picture}
\par\smallskip
\renewcommand{\thefigure}{5.3.5}
\refstepcounter{figure}
{\normalfont\bfseries 图~\thefigure}
\quad\textup{对应 (5.3.9.2)}\label{fig:5-3-5}
\end{figure}

\begin{figure}[htbp]
\centering
\setlength{\unitlength}{0.75mm}
\begin{picture}(150,98)
 \qbezier(5,54)(36,76)(67,54)
 \qbezier(5,54)(36,32)(67,54)
 \qbezier(83,54)(114,76)(145,54)
 \qbezier(83,54)(114,32)(145,54)
 \put(35,70){\circle{12}}
 \put(113,70){\circle{12}}
 \put(35,64){\circle*{1}}
 \put(113,64){\circle*{1}}
 \put(23,35){\circle{18}}
 \put(23,20){\circle{11}}
 \put(50,35){\circle{18}}
 \put(101,35){\circle{18}}
 \put(101,20){\circle{11}}
 \put(127,35){\circle{18}}
 \put(23,44){\circle*{1}}
 \put(50,44){\circle*{1}}
 \put(101,44){\circle*{1}}
 \put(127,44){\circle*{1}}
 \put(62,63){$L^{(0)}$}
 \put(62,43){$L^{(1)}$}
 \put(72,54){\makebox(0,0){$\Longrightarrow$}}
 \put(127,35){\makebox(0,0){$\scriptscriptstyle\rho_{(0),i,\alpha_0}$}}
 \put(127,47){\makebox(0,0){$\scriptstyle\tau_i^{(0)}$}}
 \put(80,8){$\scriptstyle
 \rho_{(0),i,\alpha_0}=\chi(\tau_i^{(0)})$}
\end{picture}
\par\smallskip
\renewcommand{\thefigure}{5.3.6}
\refstepcounter{figure}
{\normalfont\bfseries 图~\thefigure}
\quad\textup{对应 (5.3.9.3)}\label{fig:5-3-6}
\end{figure}

\begin{figure}[htbp]
\centering
\setlength{\unitlength}{0.8mm}
\begin{picture}(119,70)(0,-12)
 \qbezier(2,17)(24,42)(45,17)
 \qbezier(2,17)(24,-8)(45,17)
 \put(27,36){\circle{13}}
 \put(27,46.5){\circle{8}}
 \put(15,2){\circle{8}}
 \put(35,2){\circle{9}}
 \put(33,43){$\scriptstyle
 \rho_{1,i,\alpha,(j)}$}
 \put(49,22){$L^{(0)}$}
 \put(49,10.5){$L^{(1)}$}
 \put(48,17){$\Longrightarrow$}
 \qbezier(65,17)(88,42)(109,17)
 \qbezier(65,17)(88,-8)(109,17)
 \put(90,34.5){\circle{11}}
 \put(90,44){\circle{8}}
 \put(75,2){\circle{8}}
 \put(101,2){\circle{9}}
 \put(95,37){$\scriptstyle\rho_{1,i,\alpha'}$}
 \put(95,26){$\scriptstyle\rho_{1,i,\alpha}$}
 \put(112,21.5){$L^{(0)}$}
 \put(112,10.5){$L^{(1)}$}
 \put(34,-5){$\scriptstyle
 \rho_{1,i,\alpha,(j)}
 \Longrightarrow
 \rho_{1,i,\alpha}=\rho_{1,i,\alpha'}$}
 \put(50,-9.3){$\scriptstyle j\to\infty$}
\end{picture}
\par\smallskip
\renewcommand{\thefigure}{5.3.7}
\refstepcounter{figure}
{\normalfont\bfseries 图~\thefigure}
\quad\textup{对应 (5.3.9.5.1)}\label{fig:5-3-7}
\end{figure}

\begin{figure}[htbp]
\centering
\setlength{\unitlength}{0.9mm}
\begin{picture}(100,85)(0,-22)
 \qbezier(5,16)(50,47)(95,16)
 \qbezier(5,16)(50,-15)(95,16)
 \put(56,39.5){\circle{16}}
 \put(56,54){\circle{13}}
 \put(25,0){\circle{11}}
 \put(25,-9){\circle{7}}
 \put(72,-2.5){\circle{13}}
 \put(62,39){$\scriptstyle\rho_{1,i,\alpha',(j)}$}
 \put(62,27){$\scriptstyle\rho_{1,i,\alpha,(j)}$}
 \put(50,-16){\makebox(0,0){$\scriptstyle
 \lim_{j\to\infty}\rho_{1,i,\alpha,(j)}
 =
 \lim_{j\to\infty}\rho_{1,i,\alpha',(j)}$}}
\end{picture}
\par\smallskip
\renewcommand{\thefigure}{5.3.8}
\refstepcounter{figure}
{\normalfont\bfseries 图~\thefigure}
\quad\textup{对应 (5.3.9.6.1)}\label{fig:5-3-8}
\end{figure}

\begin{figure}[htbp]
\centering
\setlength{\unitlength}{0.75mm}
\begin{picture}(150,137)
 \qbezier(8,65)(75,101)(142,65)
 \qbezier(8,65)(75,29)(142,65)
 \put(40,93){\circle{26}}
 \put(22,105){\circle{14}}
 \put(105,94){\circle{28}}
 \put(105,115){\circle{14}}
 \put(105,128){\circle{11}}
 \put(42,38){\circle{27}}
 \put(42,15){\circle{18}}
 \put(104,37){\circle{27}}
 \put(40,80){\circle*{1}}
 \put(105,80){\circle*{1}}
 \put(42,51){\circle*{1}}
 \put(104,51){\circle*{1}}
 \put(22,105){\makebox(0,0){$\scriptscriptstyle\rho\ne0$}}
 \put(40,93){\makebox(0,0){$\scriptscriptstyle\rho\ne0$}}
 \put(105,94){\makebox(0,0){$\scriptscriptstyle\rho\ne0$}}
 \put(105,115){\makebox(0,0){$\scriptscriptstyle\rho\ne0$}}
 \put(105,128){\makebox(0,0){$\scriptscriptstyle\rho\ne0$}}
 \put(42,38){\makebox(0,0){$\scriptscriptstyle\rho\ne0$}}
 \put(42,15){\makebox(0,0){$\scriptscriptstyle\rho=0$}}
 \put(104,37){\makebox(0,0){$\scriptscriptstyle\rho\ne0$}}
 \put(67,83){\makebox(0,0){$\times$}}
 \put(91,94){\makebox(0,0){$\times$}}
 \put(119,94){\makebox(0,0){$\times$}}
 \put(123,74){\makebox(0,0){$\times$}}
 \put(20,58){\makebox(0,0){$\times$}}
 \put(28,38){\makebox(0,0){$\times$}}
 \put(33,15){\makebox(0,0){$\times$}}
 \put(51,15){\makebox(0,0){$\times$}}
 \put(56,38){\makebox(0,0){$\times$}}
 \put(67,47){\makebox(0,0){$\times$}}
 \put(122,58){\makebox(0,0){$\times$}}
 \put(4,70){$p$}
 \put(67,89){\makebox(0,0){$\scriptstyle P_4^{(1)}$}}
 \put(80,99){\makebox(0,0){$\scriptstyle P_3^{(1)}$}}
 \put(130,99){\makebox(0,0){$\scriptstyle P_2^{(1)}$}}
 \put(131,79){\makebox(0,0){$\scriptstyle P_1^{(1)}$}}
 \put(15,52){\makebox(0,0){$\scriptstyle P_1^{(0)}$}}
 \put(20,43){\makebox(0,0){$\scriptstyle P_2^{(0)}$}}
 \put(25,10){\makebox(0,0){$\scriptstyle P_3^{(0)}$}}
 \put(60,10){\makebox(0,0){$\scriptstyle P_4^{(0)}$}}
 \put(61,31){\makebox(0,0){$\scriptstyle P_5^{(0)}$}}
 \put(78,47){\makebox(0,0){$\scriptstyle P_6^{(0)}$}}
 \put(130,52){\makebox(0,0){$\scriptstyle P_7^{(0)}$}}
\end{picture}
\par\smallskip
\renewcommand{\thefigure}{5.3.9}
\refstepcounter{figure}
{\normalfont\bfseries 图~\thefigure}
\quad\textup{对应 (5.3.9.7.0)}\label{fig:5-3-9}
\end{figure}

\begin{figure}[htbp]
\centering
\setlength{\unitlength}{0.8mm}
\begin{picture}(115,95)(0,-10)
 \qbezier(8,50)(58,80)(108,50)
 \qbezier(8,50)(58,20)(108,50)
 \put(65,71){\circle{13}}
 \put(65,64){\circle*{1}}
 \put(35,27){\circle{20}}
 \put(35,10){\circle{12}}
 \put(65,24){\circle{20}}
 \put(35,37){\circle*{1}}
 \put(65,34){\circle*{1}}
 \put(65,24){\makebox(0,0){$\scriptscriptstyle
 \rho_{(0),i,\alpha_0}$}}
 \put(65,39){\makebox(0,0){$\scriptstyle
 \tau_i^{(0)}$}}
 \put(58,-5){\makebox(0,0){$\scriptstyle
 \chi(\tau_i^{(0)})
 =\rho_{(0),i,\alpha_0}$}}
\end{picture}
\par\smallskip
\renewcommand{\thefigure}{5.3.10}
\refstepcounter{figure}
{\normalfont\bfseries 图~\thefigure}
\quad\textup{对应 (5.3.9.8.0)}\label{fig:5-3-10}
\end{figure}

与第 4.6 节一样，
(5.3.9.5.$p$)
和 (5.3.9.6.$p$) 相互抵消
（见图 5.3.7、5.3.8）。
不难看出，
(5.3.9.3) 与 (5.3.9.8.0)
相互抵消，
(5.3.9.4) 与 (5.3.9.8.1)
相互抵消
（见图 5.3.6、5.3.10）。

对应于 (5.3.9.1) 的边界，
取纤维积后成为
（见图 5.3.4）
\[
\begin{aligned}
\phi\bigl(
 &P_1^{(1)},\ldots,P_{m'}^{(1)},
 \mathfrak n(
   P_{m'+1}^{(1)},\ldots,P_{k_1}^{(1)},
   [\ell_p,w],
   P_1^{(0)},\ldots,P_m^{(0)}),
\\[-1mm]
 &P_{m+1}^{(0)},\ldots,P_{k_0}^{(0)}
\bigr).
\end{aligned}
\]
对应于 (5.3.9.2) 的边界成为
（见图 5.3.5）
\[
\begin{aligned}
\mathfrak n\bigl(
 &P_1^{(1)},\ldots,P_{m'}^{(1)},
 \phi(
   P_{m'+1}^{(1)},\ldots,P_{k_1}^{(1)},
   [\ell_p,w],
   P_1^{(0)},\ldots,P_m^{(0)}),
\\[-1mm]
 &P_{m+1}^{(0)},\ldots,P_{k_0}^{(0)}
\bigr).
\end{aligned}
\]
对应于 (5.3.9.7.0) 的边界成为
\[
\phi\bigl(
 P_1^{(1)},\ldots,P_{k_1}^{(1)},
 [\ell_p,w],
 \widehat d(P_1^{(0)},\ldots,P_{k_0}^{(0)})
\bigr).
\]
例如，图 5.3.9 对应于
\[
\begin{aligned}
\phi_{4,6}\bigl(
 &P_1^{(1)},P_2^{(1)},P_3^{(1)},P_4^{(1)},
 [p,w],P_1^{(0)},P_2^{(0)},
\\[-1mm]
 &\mathfrak m_2(P_3^{(0)},P_4^{(0)}),
 P_5^{(0)},P_6^{(0)},P_7^{(0)}
\bigr).
\end{aligned}
\]
对应于 (5.3.9.7.1) 的边界成为\translatornote{原书实体页 314 在此式中
将首个输入写成 $P_q^{(1)}$，
但 $q$ 是输出交点，
并非左侧代数输入的下标。
译文依输入序列及上一式的对称形式
改为 $P_1^{(1)}$；
见疑误 A-0298。}
\[
\phi\bigl(
 \widehat d(P_1^{(1)},\ldots,P_{k_1}^{(1)}),
 [\ell_p,w],
 P_1^{(0)},\ldots,P_{k_0}^{(0)}
\bigr).
\]
由这些观察即得引理 5.3.8。
\end{proof}

\begin{lemma}{5.3.10}\label{lem:5-3-10}
引理 5.3.8 中的滤过 $A_\infty$
双模同态是同伦等价。
\end{lemma}
\begin{proof}
将系数约化到 $\mathbb Q$ 后，
滤过 $A_\infty$ 双模
$C(L^{(1)},L^{(0)};
  \Lambda_{0,\mathrm{nov}})$
的 $\overline{\mathfrak n}_{0,0}$
上同调是一个 $\mathbb Q$-向量空间，
其基与交点
$L^{(1)}\cap L^{(0)}$
一一对应。
另一方面，
\[
L^{(0)\prime}\cap L^{(1)\prime}
=\psi(L^{(1)}\cap L^{(0)}).
\]
因此，约化到 $\mathbb Q$ 后，
$\phi$ 诱导
$\overline{\mathfrak n}_{0,0}$
上同调的同构。
引理遂由定理 5.2.35 得出。
\end{proof}

定理 5.3.1 由引理 5.3.10
及命题 5.2.43 得出。
\end{proof}

定理 5.3.1 中构造的同伦等价，
模同伦只依赖于 $\psi$ 的辛同痕类，
与所涉及的其余选择无关。
此外，辛微分同胚的复合
对应于同伦等价的复合。
这给出下述定理 5.3.11、5.3.12。

对辛微分同胚
$\psi\colon M\to M'$，
记定理 5.3.1 中立于
$(\mathfrak g_1,\mathfrak g_0)$ 之上的
同伦等价为
\[
\varphi^\psi\colon
C(L^{(1)},L^{(0)};\Lambda_{0,\mathrm{nov}})
\longrightarrow
C(\psi(L^{(1)}),\psi(L^{(0)});
  \Lambda_{0,\mathrm{nov}}).
\]
记引理 5.3.10 中立于
$(\mathfrak f_1,\mathfrak f_0)$ 之上的
同伦等价为
\[
\phi^\psi\colon
C(L^{(1)},L^{(0)};\Lambda_{0,\mathrm{nov}})
\longrightarrow
C(\psi(L^{(1)}),\psi(L^{(0)});
  \Lambda_{0,\mathrm{nov}}).
\]
严格说来，定义
$\varphi^\psi,\phi^\psi$
还需要固定其他各种选择，
特别是多值截面 $\mathfrak s$；
因此也写作
$\varphi^{\psi,\mathfrak s}$、
$\phi^{\psi,\mathfrak s}$。

设
$\{\psi_s\}_{s\in[0,1]}$
为光滑的辛微分同胚族，
且
\[
L^{(0)\prime}=\psi_s(L^{(0)}),
\qquad
L^{(1)\prime}=\psi_s(L^{(1)})
\]
均不依赖于 $s$。

\begin{theorem}{5.3.11}\label{thm:5-3-11}
定理 5.3.1 中的 $\varphi$
模同伦不依赖于
$\{J_{\tau,t}\}_{\tau,t}$、
多值截面等选择，
而只依赖于 $\psi$ 的辛同痕类。
换言之，
$\varphi^{\psi_0,\mathfrak s_0}$
与
$\varphi^{\psi_1,\mathfrak s_1}$
同伦。
\end{theorem}

再令
$\psi'\colon M'\to M''$
为另一辛微分同胚，
并相应定义
$\varphi^{\psi'},\phi^{\psi'}$
以及
$\varphi^{\psi'\circ\psi},
 \phi^{\psi'\circ\psi}$。
\begin{theorem}{5.3.12}\label{thm:5-3-12}
$\varphi^{\psi'\circ\psi}$
与复合
$\varphi^{\psi'}\circ\varphi^\psi$
同伦。
\end{theorem}

定理 5.3.11、5.3.12 的证明
与第 5.3.3、5.3.4 小节
将分别给出的定理 5.3.27、5.3.39
的证明类似而且更容易，
故此处从略。
